\documentclass[11pt]{article}
\usepackage[margin=1in]{geometry}
\usepackage{amsmath,amssymb}
\usepackage{setspace}
\setstretch{1.2}

\title{Referee Report: ``Multiproduct Intermediaries'' \\ \small Rhodes, Watanabe, and Zhou (2021)}
\author{}
\date{}

\begin{document}
\maketitle
\vspace{-2em}

\section*{Summary}

This paper asks why multiproduct intermediaries are profitable even when they do not reduce search costs or offer lower prices. The answer is assortment choice: by stocking high consumer-surplus products exclusively, the intermediary attracts visits, and then profits from selling other high-margin products to consumers who are already in the store. The paper characterizes the optimal product range and applies the framework to shopping malls and direct-to-consumer (DTC) sales.

\section*{Model}

A continuum of manufacturers (measure one) each produce a different product $i$ with marginal cost $c_i \geq 0$; consumer demand for product $i$ at price $p_i$ is $Q_i(p_i)$. A unit mass of consumers have identical preferences but differ in their search cost $s \sim F$ on $[0, \bar{s}]$: a consumer must pay $s$ to visit a seller and learn its price. Each product is summarized by two statistics -- monopoly profit and consumer surplus:
\[
\pi_i = (p_i^m - c_i)\, Q_i(p_i^m), \qquad v_i = \int_{p_i^m}^{\infty} Q_i(p)\, dp,
\]
where $p_i^m$ is the monopoly price. Since consumers do not observe prices before searching, all sellers charge $p_i^m$ (Diamond, 1971). The intermediary makes take-it-or-leave-it two-part tariff offers to manufacturers with wholesale price $\tau_i = c_i$ and lump-sum $T_i = \pi_i F(v_i)$ (the manufacturer's outside option from direct sales), and can request exclusivity. The key reduction: each product is fully described by $(\pi_i, v_i)$ in a product space $\Omega \subset \mathbb{R}_+^2$ with joint distribution $G$.

\paragraph{Simple case.} Assume all contracts are exclusive, $h(m) = m$ (where $m = \int_A dG$ is the measure of products stocked), and no capacity constraint ($\bar{m} = 1$). Let $A \subseteq \Omega$ be the stocked set. A consumer of type $s$ visits the intermediary iff $s \leq \hat{v}$, where $\hat{v} = \int_A v\, dG / \int_A dG$ is the average consumer surplus across stocked products. The net profit from stocking $(\pi, v)$ is $\pi[F(\hat{v}) - F(v)]$: the intermediary attracts $F(\hat{v})$ consumers but pays the manufacturer $\pi F(v)$. Products with $v < \hat{v}$ are profitable; those with $v > \hat{v}$ generate a loss. With stocking indicator $q(\pi,v) \in \{0,1\}$:
\[
\max_{q} \int_\Omega q(\pi,v)\, \pi [F(\hat{v}) - F(v)]\, dG \qquad \text{s.t.} \quad \hat{v} = \frac{\int_\Omega q\, v\, dG}{\int_\Omega q\, dG}.
\]
Using a Lagrange multiplier $\lambda > 0$, the value of stocking $(\pi, v)$ decomposes into a \emph{direct effect} $\pi[F(\hat{v}) - F(v)]$ and an \emph{indirect effect} $\lambda(v - \hat{v})$, capturing how the product changes average surplus and thus the visiting mass. These have opposite signs. The optimal assortment has two negatively correlated regions: high-$v$, low-$\pi$ (traffic generators) and low-$v$, high-$\pi$ (profit generators).

\paragraph{General case.} Allow non-exclusive contracts, a general search technology $h(m)$ (cost of visiting the intermediary is $h(m) \times s$ when it stocks $m$ products), and a capacity constraint $m \leq \bar{m}$. The core trade-off carries over. High-$v$ products are stocked exclusively to attract consumers; low-$v$, high-$\pi$ products generate profit. With economies of search ($h'(m) < 1$), all high-$v$ products are stocked; with diseconomies ($h'(m) > 1$), only the low-$\pi$ ones among them. Larger capacity leads to more products stocked but a smaller exclusive share.

\paragraph{The fundamental assumption.} The mechanism rests on an asymmetry in search costs. Buying from manufacturers costs $s$ \emph{per firm} visited ($n$ firms $= n \times s$), but visiting the intermediary costs a \emph{single} $h(m) \times s$ that gives access to all $m$ products. Even when $h(m) = m$, the consumer makes one visit decision that opens the full assortment. This is what lets the intermediary bundle traffic generators with profit generators: a consumer who would never search for a low-$v$, high-$\pi$ product on its own will buy it once already inside the store. Without this one-stop-shopping structure, cross-product externalities vanish and the intermediary cannot profitably exist.

\section*{Applications}

\textbf{Shopping malls.} The mall is a platform that does not set prices itself. High-$v$ stores play the role of ``anchor stores'' and are subsidized to attract foot traffic; low-$v$, high-$\pi$ stores pay rent. The $(\pi, v)$ framework applies directly.

\textbf{DTC expansion.} When direct sales become easier (parameterized by $\theta > 0$ reducing the consumer's search cost for manufacturers), the manufacturer's outside option rises to $\pi F(v/\theta)$. The intermediary responds by carrying fewer products, stocking a larger share of them exclusively, and its selection becomes more polarized in $(\pi, v)$ space. Its profit falls.

\section*{Assessment}

The $(\pi, v)$ reduction is the main contribution. It turns a complicated combinatorial assortment problem into a tractable two-dimensional one, and the direct/indirect effect decomposition gives clear intuition for why the intermediary stocks what it stocks. The applications to malls and DTC are natural and informative.

Several points deserve discussion:

\textit{Monopoly pricing.} The whole analysis rests on Diamond (1971): consumers do not observe prices before search, so every seller charges the monopoly price $p_i^m$. This rules out price competition entirely. The authors argue what matters is that all sellers of a given product charge the same price, not the level itself. Fair enough, but when the intermediary stocks a product non-exclusively it does compete with the manufacturer in reality, and prices should react. Checking robustness to price responses would strengthen the paper.

\textit{Single intermediary.} There is one intermediary with all bargaining power, so it fully extracts the traffic-generation rents. With competing intermediaries, high-$v$ manufacturers would have outside options beyond direct sales -- they could threaten to join a rival. This should erode intermediary profits and change assortment composition. The paper acknowledges this as future work and notes the $(\pi, v)$ approach may not extend directly. I sketch one direction below.

\textit{Consumer heterogeneity.} Consumers differ only in $s$, not in tastes. This is needed for the $(\pi,v)$ representation since it guarantees all sellers charge the same price. Section 7.2 allows some demand heterogeneity and the main insights survive, but only under restrictions that prevent differential pricing across channels. Richer heterogeneity would likely affect both stocking and contracting decisions.

\textit{The one-stop-shopping assumption.} As noted above, the mechanism depends on consumers paying a single cost to access the intermediary's full assortment, rather than a per-product cost as with manufacturers. This is a strong assumption. In some retail settings (e.g.\ online platforms where consumers still need to search within the platform), the cost structure may be closer to per-product. The paper could discuss more explicitly how results change if $h(m)$ grows faster than linearly, making in-store search increasingly costly.

\section*{Extension Sketch: Competing Platforms}

Consider two symmetric platforms $j \in \{1, 2\}$ bidding simultaneously for products via exclusive contracts. Let $A_j \subseteq \Omega$ be platform $j$'s exclusive set ($A_1 \cap A_2 = \emptyset$), with $\hat{v}_j = \int_{A_j} v\, dG / \int_{A_j} dG$. Platform $j$'s profit from $(\pi, v) \in A_j$ is $\pi[F(\hat{v}_j) - F(v)]$. The difference from the baseline: platforms bid against each other for high-$v$ products. The indirect value of such a product to platform $j$ is roughly $\lambda_j(v - \hat{v}_j)$; competition pushes the transfer up to the rival's willingness to pay, so the manufacturer captures approximately $T^* \approx \pi F(v) + \lambda_{-j}(v - \hat{v}_{-j})$ -- its direct-sales outside option plus the indirect value it would create for the competitor.

\textbf{What likely survives:} the negatively correlated stocking pattern in $(\pi, v)$, since the traffic/profit trade-off persists; the direct/indirect decomposition within each platform.

\textbf{What breaks:} the $(\pi, v)$ representation gets harder, because a consumer now chooses \emph{between} two intermediaries. The threshold $\hat{v}_j$ depends on $A_{-j}$ -- a consumer may skip platform $j$ not because $s$ is high, but because the rival offers better surplus. This interdependence makes the stocking problem non-separable across products. Existence of a pure-strategy equilibrium in assortments is not obvious.

\textbf{Conjecture:} each platform stocks fewer high-$v$ exclusives (they cost more to acquire), platform profits fall, and welfare likely increases as competition curbs excessive exclusivity.

\end{document}
